\documentclass{ieeeaccess}
\usepackage{cite}
\usepackage{amsmath,amssymb,amsfonts}
\usepackage{algorithmic}
\usepackage{graphicx}
\usepackage{textcomp}
\usepackage{longtable}
\usepackage{multirow}
\usepackage{array}
%\usepackage{tikz}
%\usetikzlibrary{fit,arrows,shapes,positioning}
%\usepackage{pgfplots}
%\pgfplotsset{compat=1.18}
%\usepgfplotslibrary{groupplots}
%\usepackage{mathptmx}   % or newtxtext,newtxmath
\usepackage{booktabs}
%\usepackage{xcolor}

%\usepackage[T1]{fontenc}
%\usepackage[margin=1in]{geometry}
\usepackage{newtxtext,newtxmath}
\usepackage[table]{xcolor}
%\usepackage{array}
\usepackage{booktabs}

\definecolor{cellred}{HTML}{E74C3C}
\definecolor{low}{HTML}{F47C40}
\definecolor{mid}{HTML}{F6D365}
\definecolor{highmid}{HTML}{B7E08A}
\definecolor{green88}{HTML}{A5D683}
\definecolor{green89}{HTML}{94CC7B}
\definecolor{green90}{HTML}{82C274}
\definecolor{green92}{HTML}{4E8B5A}
\definecolor{high}{HTML}{2D6A4F}
\definecolor{peak}{HTML}{1A4D38}
\definecolor{headergray}{HTML}{F2F2F2}
\definecolor{labelgray}{HTML}{F7F7F7}
\definecolor{gridgray}{HTML}{C4C4C4}

\newcommand{\best}[1]{\textbf{#1}}
\newcommand{\redcell}[1]{\cellcolor{cellred!55}#1}
\newcommand{\lowcell}[1]{\cellcolor{low!55}#1}
\newcommand{\midcell}[1]{\cellcolor{mid!60}#1}
\newcommand{\highmidcell}[1]{\cellcolor{highmid!60}#1}
\newcommand{\greenAcell}[1]{\cellcolor{green88!65}#1}
\newcommand{\greenBcell}[1]{\cellcolor{green89!65}#1}
\newcommand{\vhighcell}[1]{\cellcolor{green90!65}#1}
\newcommand{\topcell}[1]{\cellcolor{green92!65}#1}
\newcommand{\highcell}[1]{\cellcolor{high!45}#1}
\newcommand{\peakcell}[1]{\cellcolor{peak!60}#1}
\newcommand{\modelname}[1]{\cellcolor{labelgray}\textbf{#1}}
\newcommand{\groupgap}{\noalign{\vskip 0.24em}}
\newcommand{\splitline}{\cline{1-6}\cline{8-12}}
% FPR95 label and color commands for dual-row tables (lower FPR95 = better = darker green)
\newcommand{\fprname}{\cellcolor{labelgray}\scriptsize\hspace{0.3em}FPR95\,$\downarrow$}
\newcommand{\fprpeakcell}[1]{\cellcolor{peak!60}#1}
\newcommand{\fprhighcell}[1]{\cellcolor{high!45}#1}
\newcommand{\fprtopcell}[1]{\cellcolor{green92!65}#1}
\newcommand{\fprvhighcell}[1]{\cellcolor{green90!65}#1}
\newcommand{\fprmidcell}[1]{\cellcolor{mid!60}#1}
\newcommand{\fprlowcell}[1]{\cellcolor{low!55}#1}
\newcommand{\fprredcell}[1]{\cellcolor{cellred!55}#1}

\newcolumntype{L}{>{\raggedright\arraybackslash}m{3.5cm}}
\newcolumntype{C}{>{\centering\arraybackslash}m{0.9cm}}
\newcolumntype{G}{>{\centering\arraybackslash}m{0.15cm}}

\usepackage{bm}
\makeatletter
\AtBeginDocument{\DeclareMathVersion{bold}
\SetSymbolFont{operators}{bold}{T1}{times}{b}{n}
\SetSymbolFont{NewLetters}{bold}{T1}{times}{b}{it}
\SetMathAlphabet{\mathrm}{bold}{T1}{times}{b}{n}
\SetMathAlphabet{\mathit}{bold}{T1}{times}{b}{it}
\SetMathAlphabet{\mathbf}{bold}{T1}{times}{b}{n}
\SetMathAlphabet{\mathtt}{bold}{OT1}{pcr}{b}{n}
\SetSymbolFont{symbols}{bold}{OMS}{cmsy}{b}{n}
\renewcommand\boldmath{\@nomath\boldmath\mathversion{bold}}}
\makeatother
\def\BibTeX{{\rm B\kern-.05em{\sc i\kern-.025em b}\kern-.08em
    T\kern-.1667em\lower.7ex\hbox{E}\kern-.125emX}}

%Your document starts from here ___________________________________________________
\begin{document}
\history{Date of publication xxxx 00, 0000, date of current version xxxx 00, 0000.}
\doi{10.1109/ACCESS.2024.0429000}

\title{Out-of-Distribution Detection with Spike-Like Networks: Population Coding, Training Depth, and Single-Step Feature Representations}
\author{Aleksej Avramovi\'{c}\authorrefmark{1,2},
Davor Sluga\authorrefmark{3}, Slavica Gaji\'{c}\authorrefmark{2}, Vedran Jovanovi\'{c}\authorrefmark{2,1},
\IEEEmembership{Member, IEEE}, and Vladimir Risojevi\'{c}\authorrefmark{2}}

\address[1]{HTEC Group, Banja Luka, Bosnia and Herzegovina}
\address[2]{Faculty of Electrical Engineering, University of Banja Luka, Bosnia and Herzegovina}
\address[3]{Faculty of Computer and Information Science, University of Ljubljana, Slovenia}
\tfootnote{This paragraph of the first footnote will contain support
information, including sponsor and financial support acknowledgment. For
example, ``This work was supported in part by the U.S. Department of
Commerce under Grant BS123456.''}

\markboth
{A. Avramovi\'{c} \headeretal: Out-of-Distribution Detection with Spike-Like Networks}
{A. Avramovi\'{c} \headeretal: Preparation of Papers for IEEE TRANSACTIONS and JOURNALS}

\corresp{Corresponding author: Aleksej Avramovi\'{c} (e-mail: aleksej.avramovic@ etf.unibl.org).}


\begin{abstract}
Out-of-distribution (OoD) detection is a critical requirement for deploying deep learning systems in open-world environments, yet most existing methods rely on dense, full-precision convolutional networks with high computational and energy demands. This paper investigates spike-like models, based on standard convolutional backbones in which conventional activations are replaced by leaky integrate-and-fire neurons, as a computationally efficient alternative for OoD detection. Constrained to a single simulation time step at inference, these models produce inherently binary spike activations while retaining the feature extraction capacity of the underlying backbone, substantially reducing arithmetic complexity relative to conventional CNNs. We systematically study the interaction between training objectives, feature representation type, population coding, and post-hoc OoD detection methods across multiple backbone architectures and in-distribution benchmarks including CIFAR-10, SVHN, and CIFAR-100. Two families of detectors are evaluated: distance-based methods operating on penultimate-layer membrane-potential features, and scoring-based post-hoc methods applied to final-layer representations. Experimental results show that spike-like models trained with a mean-squared-error objective produce penultimate-layer features that support strong $k$-nearest-neighbor OoD detection, matching or exceeding conventional CNN baselines at a fraction of the arithmetic cost. Population coding consistently improves OoD separability and is identified as a prerequisite for multi-step training to yield further gains under single-step inference. 
%Comment that the proposed approach matches or outperormes the best results given in OpenOOD. 
\end{abstract}

\begin{keywords}
Anomaly detection, convolutional neural networks, distance-based detection, energy-efficient inference, feature representations, k-nearest neighbors, leaky integrate-and-fire neurons, membrane potentials, neuromorphic computing, out-of-distribution detection, population coding, post-hoc detection, residual networks, scoring-based detection, single-step inference, spike patterns, spiking neural networks, surrogate gradient training.
\end{keywords}

\titlepgskip=-21pt

\maketitle

% TODO Keep the latex formatting, tabes and references.
% The text is more or less solid, but needs fine tuning for the following.
%TODO: harmozine therminology for: single-step spiking network, single-step SNN, hard-quantized network, signle-step spike-like network
%Suggest the unified therminology and if needed, suggest the way to define it in the introduction
% TODO Keep the passive voice. Fix spelling and grammar. Keep the fluency.

%Keep the existing single-sentence definition you added, then (a) add a 2–3 line glossary/footnote listing: "CNN = standard non‑spiking conv nets; SNN = fully spiking (multi‑step); spike‑like = single‑step SNN implemented on a CNN backbone (CNN+LIF); spike‑based = adjective for features produced by any spiking mechanism." After that, perform a small pass applying the specific replacements above.


\section{Introduction}
\label{sec:introduction}

\PARstart{M}{achine} learning (ML) and artificial intelligence (AI) have become central to modern industrial, scientific, and societal applications. From autonomous vehicles and robotics to medical diagnostics and predictive maintenance, intelligent systems are increasingly deployed in environments where reliability and adaptability are critical. The success of such systems is determined not only by performance on known data but also by robustness when unexpected or unfamiliar situations are encountered.

One of the major challenges in this context is the out-of-distribution (OoD) detection task, which concerns identifying inputs that deviate from the data distribution observed during training. In safety-critical domains, a failure to recognize such unfamiliar inputs can lead to incorrect or hazardous outcomes. As a result, OoD detection has emerged as a fundamental component of trustworthy and safe AI, ensuring that models can distinguish between familiar and novel data and respond appropriately~\cite{hendrycks2017baseline, liu2020energy, lee2018simple}.

Convolutional Neural Networks (CNNs) remain a common choice for OoD detection thanks to their strong feature extraction and mature architectures, although recent studies have also investigated transformer-based backbones for this task. A variety of post-hoc scoring methods, including Maximum Softmax Probability (MSP)~\cite{hendrycks2017baseline}, energy-based scores~\cite{liu2020energy}, and Mahalanobis-distance-based detectors~\cite{lee2018simple}, have been proposed to assess predictive confidence and identify OoD samples without retraining the model. These approaches provide strong baselines in the literature. However, CNN-based inference relies on dense floating-point activations and a large number of multiply--accumulate (MAC) operations, resulting in significant computational and energy demands that limit applicability in resource-constrained, embedded, or real-time environments.

Spiking Neural Networks (SNNs) offer a promising alternative due to their event-driven and sparse computational paradigm. In SNNs, neurons emit discrete spikes only when their membrane potentials exceed a threshold, leading to reduced activity and potentially lower energy consumption during inference. In particular, when inference is performed with a small number of time steps, and in the extreme case with a single time step, the number of synaptic operations can be significantly reduced. However, this efficiency can trade off against representational capacity: high classification accuracy in SNNs is often achieved via rate-based coding and long temporal integration windows, which require many time steps and diminish potential energy savings. Consequently, a fundamental trade-off arises between accuracy and efficiency in SNN-based systems.

Population coding has been proposed to mitigate these limitations. By representing information across groups of neurons rather than single units, population coding can improve robustness, reduce sensitivity to noise, and stabilize training dynamics. In SNNs, population coding may approximate rate-based representations even when operating with few time steps, making it attractive for low-latency and energy-efficient inference regimes where traditional temporal coding strategies are impractical.

We use 'spike-like' to denote single‑step spiking models implemented on conventional convolutional backbones (i.e., CNN architecture + LIF activations); when we write 'SNN' we mean fully spiking networks (multi-step or event-driven), and when we write 'CNN' we mean standard non-spiking convolutional models.

Motivated by given observations, we study whether spike-like models can provide competitive OoD detection when constrained to low-latency, low-compute inference regimes (i.e., reduced time-steps and arithmetic). Specifically, we address the following research questions: (i) whether training strategies can be adapted to obtain sufficiently discriminative feature representations under single-step inference, and (ii) whether standard post-hoc OoD detection methods can be effectively applied to such representations. In addition, the role of population coding in enhancing feature separability and robustness is examined. 

Related work~\cite{Olber2022DetectionOO} has shown that binary activation patterns extracted from CNNs can capture internal network responses that are informative for OoD detection. In this study, such binary representations are learned directly via spiking activations during training, enabling feature vectors that are inherently suitable for OoD detection to be obtained from scratch. A fair and direct comparison with CNN-based baselines is enabled by this design choice, while key spiking characteristics--namely binary activations, sparse computation, and reduced arithmetic complexity--are preserved, thereby allowing the trade-off between discriminative performance and computational efficiency to be systematically analyzed.

In this context, the direct training of spike-like models is investigated: the spike activations of these networks are inherently binary and therefore quantized outputs are produced without conventional quantization methods being applied. In such single-step models, feature representations suitable for OoD detection can be learned from scratch.
%, without reliance on post-hoc conversion from pre-trained CNNs.

The main contributions of this work are:
\begin{itemize}
  \item A spiking-based framework in which spiking activation dynamics are incorporated into standard convolutional backbones, enabling reduced computational complexity and energy-efficient inference.
  \item A systematic analysis of training strategies for single-step and multi-step SNNs, with particular emphasis on achieving competitive performance under single-step inference.
  \item An evaluation of standard post-hoc OoD detection methods on spike-patterns and membrane-potential-based feature representations.
  \item An investigation of the impact of population coding on feature robustness, separability, and OoD detection performance across multiple architectures and datasets.
  \item An exploration of directly trained, inherently binarized spike-like models, avoiding reliance on CNN-to-SNN conversion pipelines.
\end{itemize}

The remainder of this paper is organized as follows. Section~II reviews related work on CNN- and SNN-based OoD detection. Section~III describes the proposed architecture and training methodology. Section~IV presents the experimental setup and results. Section~V concludes the paper and outlines future research directions.



\section{Background (Research Context)}
\label{sec:background}

The OoD detection task task is to decide whether a test sample $x \in \mathcal{X}$ originates from the training distribution $P_{\text{train}}(x)$. When a sample is out-of-distribution (i.e. $x \nsim P_{\text{train}}(x)$) feature representations and confidence scores can be misleading, so detectors flag such inputs to prevent overconfident or erroneous decisions. Let $P_{\text{train}}(x, y)$ denote the joint distribution of inputs and labels observed during training, and $P_{\text{train}}(x)$ its marginal. OoD detection can then be formulated as a binary hypothesis test:
\begin{equation}
\label{eq:01a}
\mathcal{H}_0: \; x \sim P_{\text{train}}(x),
\end{equation}
\begin{equation}
\label{eq:01b}
\mathcal{H}_1: \; x \sim P_{\text{out}}(x), \quad P_{\text{out}} \neq P_{\text{train}},
\end{equation}
where $\mathcal{H}_0$ corresponds to in-distribution samples and $\mathcal{H}_1$ to out-of-distribution samples. Reliable OoD detection is essential for deploying deep learning models in open-world environments, where unexpected or anomalous samples may occur.

Conventional CNN-based inference relies on dense convolutional operations, resulting in significant computational load and power consumption during feature extraction and classification. This cost is largely independent of input activity, since all layers are evaluated exhaustively at inference time. While CNNs provide strong representational capacity and state-of-the-art OoD performance, their deployment in energy-constrained environments remains challenging.

SNN provide an alternative computational paradigm based on sparse and event-driven processing. Neurons emit spikes only when their membrane potentials exceed a threshold, allowing computation to be performed selectively. Under appropriate sparsity conditions and short inference windows, this can reduce both computational cost and energy consumption. However, these advantages depend strongly on spike activity, encoding strategies, and network configuration.

In this work, OoD detection is studied on static images, and we investigate whether temporal dynamics can be constrained (e.g., to single‑step inference) without degrading detection performance. Using single-step regime reduces arithmetic complexity while preserving key spiking characteristics such as sparse, event-driven activations.

Feature representations can be obtained either as binary spike patterns or as continuous membrane potentials from intermediate network layers. Input images are used directly as grayscale or RGB inputs rather than being converted into temporal spike sequence; the first convolutional layer together with the spike activation function implicitly converts pixel intensities into spike activations (this conversion is part of the spiking backbone not a separate CNN preprocessor), and subsequent layers operate only on spike activations. This design avoids the overhead of rate- or time-based encoding schemes and supports low-latency inference.

The network models are based on standard convolutional and residual architectures in which conventional activations are replaced with spiking leaky integrate-and-fire (LIF) neurons. Training is performed using surrogate-gradient methods, due to non-differentiability of spikes. Both single-step and multi-step training regimes are considered, while inference is restricted to a single time step.
% for the experiments focused on low-latency operation.

Within this framework, population coding is investigated as a mechanism for improving feature robustness and separability. By distributing information across groups of neurons, population coding can enhance resilience to noise and improve the discriminative structure of feature representations used for OoD detection.

\section{Related work}
\label{sec:related}

In many real-world applications, such as autonomous driving or industrial monitoring, systems must detect inputs that deviate from the training distribution in order to ensure safe and reliable operation. This requirement has motivated extensive research on OoD detection in deep learning.

\vspace{0.3em}
\noindent\textbf{CNN-based OoD detection.} 
Many OoD approaches are applied post-hoc on feature representations extracted from a trained classifier. Broadly, post-hoc methods used in the literature fall into two families. Scoring-based methods compute a scalar anomaly score from a single feature vector or logit (examples: Maximum Softmax Probability (MSP)~\cite{hendrycks2017baseline}, energy-based scores~\cite{liu2020energy}, Virtual Logit Matching (ViM)~\cite{wang2022vim}, and ASH~\cite{djurisic2022ash}). Distance-based methods derive scores from distances in feature space to in-distribution references, e.g., Mahalanobis-distance detectors~\cite{lee2018simple}, Nearest Class Mean (NCM), and $k$-Nearest Neighbors (kNN)~\cite{sun2022knn}. 

Beyond post-hoc scoring, other strategies include explicit feature-space density estimation and density-ratio models~\cite{sun2022out, zhang2022ood}, clustering-based representations~\cite{sinhamahapatra2022all}, and methods that incorporate calibration or small auxiliary training stages (e.g., ODIN-like temperature/scaling schemes). Improved evaluation protocols and benchmarks have also refined how methods are compared~\cite{kuan2022back, yang2022openood}. While most of the algorithms above were developed for convolutional networks, many can be adapted to spiking networks by using either spike patterns or membrane potentials; the efficacy and practical considerations of transferring these post-hoc techniques to sparse, binary representations remain an active area of research.

\vspace{0.3em}
\noindent\textbf{OoD detection with spiking networks.}
OoD detection in spiking neural networks has only recently been explored. In~\cite{2023-martinez}, spike-count activations across layers were used to derive explainable OoD indicators, achieving competitive performance on both static and event-based datasets. In~\cite{terresescudero2025}, the robustness of fully spiking networks in open-world scenarios was investigated, demonstrating improved generalization under forward-only training schemes.

Uncertainty estimation in SNNs has also been studied as a related problem. Spike-based confidence measures~\cite{Sun_2023} and Bayesian formulations~\cite{prabodh2024} indicate that spike activity statistics can provide meaningful uncertainty estimates that can be leveraged for OoD detection.

\vspace{0.3em}
\noindent\textbf{Training and architectural advances in SNNs.}
Recent progress in surrogate-gradient learning has enabled effective training of deep spiking networks using gradient-based optimization~\cite{neftci2019surrogate, zenke2021remarkable, eshraghian2023training}. Architectures such as Spiking-ResNet~\cite{2022-fang, 2021-zheng} and directly trained deep SNNs~\cite{2020-rathi} have achieved performance comparable to CNNs on standard benchmarks. Additional improvements include information-theoretic loss functions~\cite{2022-Guo} and learnable membrane dynamics~\cite{2025-shen}, which enhance representation quality under limited temporal depth.

\vspace{0.3em}
\noindent\textbf{Population coding in spiking networks.}  
Population coding distributes information across groups of neurons, leading to robust and noise-tolerant representations~\cite{Tkacik2010Optimal, Tripp2012Population}. Its application in machine learning includes temporal classification~\cite{Pan2019NeuralPopulation} and reinforcement learning~\cite{Zhang2021PopulationCoding}. However, its role in OoD detection has not been extensively studied. In this work, population coding is evaluated as a mechanism for improving feature separability in both spiking and dense representations.

\vspace{0.3em}
\noindent\textbf{Computational and energy considerations.}  
The computa-tional complexity of CNNs is typically measured using floating-point operations (FLOPs) or multiply-accumulate (MAC) counts. In contrast, SNNs rely on spike-triggered synaptic operations (SynOps), which depend on activity and sparsity~\cite{roy2019towards, davies2018loihi, plagwitz2023to}. For spike-like models, deeper layers primarily involve accumulate-only (AC) operations, reducing arithmetic complexity~\cite{chen2023training}. Comparisons between MACs and SynOps are commonly used as a hardware-agnostic proxy for computational efficiency~\cite{lemaire2023estimation}.

The key characteristics of CNNs, SNNs (multi‑step fully spiking), and spike-like models (single‑step CNN with LIF) are summarized in Table~\ref{tab:energy-comparison-summary}.

\begin{table*}[tbp]
\caption{Comparison of computational and energy characteristics for CNNs, SNNs and spike-like models. Energy estimates are expressed qualitatively, based on existing literature.}
\begin{center}
{\renewcommand{\arraystretch}{0.95}\small
{\small
\begin{tabular}{l >{\raggedright\arraybackslash}m{3.1cm} >{\raggedright\arraybackslash}m{3.1cm} >{\raggedright\arraybackslash}m{3.1cm} >{\raggedright\arraybackslash}m{3.1cm}}
\toprule
\textbf{Network Type} & \textbf{Dominant operations} & \textbf{Operation count metric} & \textbf{Energy proxy} & \textbf{Sparsity / event-driven} \\
\midrule
CNN & Multiply-and-accumulate operations in all layers & FLOPs / MACs & High: dense floating-point operations & Low: all layers evaluated for every input \\
\midrule
SNN & Synaptic accumulations triggered by spikes & SynOps / spike counts & Potentially low: depends on spike activity and sparsity & High: computation occurs only when spikes are generated \\
\midrule
Spike-like model & First layer: MACs; deeper layers: ACs on binary activations & MACs and ACs & Reduced: lower arithmetic complexity while preserving spiking properties & Moderate: sparse activations with single-step inference \\
\bottomrule
\end{tabular}
}
}
\label{tab:energy-comparison-summary}
\end{center}
\end{table*}


\section{Network Architectures}
\label{sec:network}

The implemented network architectures are based on convolutional backbones augmented with key properties of spiking neural networks. Standard ReLU activations used in CNNs are replaced by LIF neurons, which produce binary spike outputs while preserving spatial feature extraction capabilities. Both sparse feature vectors and dense feature representations can be extracted and subsequently used for post-hoc OoD detection. These representations can be obtained either as sparse spike patterns (1 indicating a spike and 0 indicating silence) or as dense feature vectors derived from membrane potentials in the penultimate or output layer.

The number of simulation time steps employed during training and inference can be selected flexibly. In prior work on OoD detection with spiking networks, inference commonly required on the order of tens of time steps; studies emphasizing classification accuracy sometimes operated with fewer steps.

\begin{table*}[tbp]
\caption{Summary of network backbones. All models use LIF activations ($\beta=0.95$, threshold $=0.25$, reset $=0$), surrogate-gradient training, and optional population coding in the output layer. The number of parameters depends on the input image size and the number of classes.}
\begin{center}
\small
\begin{tabular}{l l p{10cm}}
\toprule
\textbf{Backbone} & \textbf{Depth} & \textbf{Key Properties} \\
\midrule
spike-ConvNet & 2 conv + 3 FC & Conv($3{\times}3$)~$\rightarrow$~LIF~$\rightarrow$~AvgPool~$\rightarrow$~Conv($3{\times}3$)~$\rightarrow$~LIF~$\rightarrow$~FC(500, 300, $C$)~$\rightarrow$~LIF \\
\midrule
spike-ResNet10 & 10 conv + 1 FC & Residual blocks [64, 128, 256, 512], conv2d+BN+LIF, one block per stage, skip connections with binary clamping \\
\midrule
spike-ResNet18 & 18 conv + 1 FC & Residual blocks [64, 128, 256, 512], conv2d+BN+LIF, two blocks per stage, selected super-blocks with four convolutional layers, skip connections clamped to 1 \\
\bottomrule
\end{tabular}
\label{tab:network-backbones}
\end{center}
\end{table*}

In the limiting case, the number of simulation time steps can be reduced to one. Under this configuration, the network operates as a binarized spike-like model: activations are binary and the temporal dynamics of LIF neurons are largely suppressed, resulting in integrate-and-fire (IF) behavior. A neuron emits a spike only if sufficient input activity from the previous layer causes its membrane potential to exceed the firing threshold. Although explicit temporal processing is not exploited, this single-step regime preserves the sparsity and computational efficiency associated with spiking networks while enabling the spiking-augmented convolutional backbone to produce discriminative feature representations: spike patterns and and continuous membrane-potential vectors. Both representation types are suitable for OoD detection.

The architecture can be shifted toward a more temporally rich spiking regime by increasing the number of simulation time steps during training, thereby enabling the temporal dynamics of LIF neurons to be exploited. To preserve low inference-time power consumption, networks trained with multiple time steps can still be used with single-step propagation during feature extraction.

In this work, given backbone architectures are combined with spiking activation layers. Depending on the underlying backbone, the resulting models are referred to as Spike-ConvNet or Spike-ResNet architectures. In all cases, LIF activations are employed with parameters $\beta$ equal to 0.95, threshold 0.25, and zero reset. Residual blocks in Spike-ResNet architectures consist of multiple convolutional layers followed by BatchNorm2d and LIF activations. Skip connections are applied to spiking feature maps using element-wise addition, after which outputs are clamped to a maximum value of 1. Downsampling is performed using convolutional layers with stride 2 followed by LIF activation, ensuring that spiking characteristics are preserved across network stages. The key architectural characteristics of all implemented backbones are summarized in Table~\ref{tab:network-backbones}.

Population coding is implemented by expanding the final linear layer to size $C\times E$ (where $C$ is number of output classes and  $E$ is the expansion factor): the neurons from the penultimate layer are fully connected to all $C\times E$ neuron from the final layer. The per-class response is then computed by aggregating $E$ outputs for each class, which only increases final-layer parameters while leaving the convolutional backbone unchanged.


\section{Methodology}
\label{sec:methodology}

\subsection{Population coding}
\label{sec:population-coding}

Population coding is optionally applied at the output layer by assigning each class to a population of output neurons rather than a single output unit. For a given expansion value, the final-layer output vector is partitioned into class-wise populations, and a class-level response is obtained by aggregating the activity within each population. The aggregation rule is matched to the loss function: summed spike counts are used for spike-count objectives, whereas mean-valued membrane potentials are used for membrane-based objectives. The loss is then applied to these aggregated class-level responses.

Through this mechanism, training is configured so that the population corresponding to the target class exhibits stronger collective activity, while populations associated with other classes remain suppressed. This distributed representation can reshape the learned feature space, potentially improving intra-class compactness and inter-class separation.


\subsection{Training Procedure}
All backbones were trained under identical experimental conditions. The loss functions, preprocessing and augmen-tation choices, population-coding setup, and optimization details used in the experiments are summarized below. The first three approaches are also used for spike-like models operating with a single simulation time step ($T=1$). In this regime, the networks retain hard-quantized spiking activations while effectively operating in a single-step inference mode.

\begin{itemize}
\item \textbf{SNN + MSE (spike-pattern-based mean-square error):} Optimization is performed on the final-layer spike pattern. Spike counts are accumulated over the simulation window, and the mean-squared error with respect to the target representation is minimized. No per-channel normalization is applied for this objective, and data augmentation is restricted to geometric transformations (random horizontal flip and random crop).

\item \textbf{SNN + CL (spike-count-based):} Optimization is performed directly on integer spike-count targets. As with MSE-based training, no per-channel normalization is applied, and augmentation is restricted to geometric transformations.

\item \textbf{SNN + CE (membrane-potential-based):} Optimization is performed on the continuous membrane potentials of the final layer. Standard per-channel mean--standard deviation normalization is applied, together with both geometric and photometric augmentations.

\item \textbf{CNN + CE:} Identical backbones are trained as regular convolutional neural networks using cross-entropy loss on output logits to provide a conventional CNN baseline for comparison. Standard normalization and both geometric and photometric augmentations are applied.
\end{itemize}

When population coding is used, final-layer outputs are reshaped into class-wise populations and aggregated before applying the loss (see Subsection~\ref{sec:population-coding}).
Optimization was carried out using the Adam optimizer with standard parameters and a weight decay of $10^{-4}$. A cyclical learning-rate schedule (OneCycleLR) with an initial learning rate of $2\cdot10^{-4}$ was employed to improve convergence over 400 training epochs. All models were trained with a batch size of 64.

\subsection{Feature Extraction for OoD Detection}
Feature vectors are extracted from either the penultimate or final layers and may consist of binary spike patterns or continuous membrane potentials depending on the evaluation protocol. All OoD methods considered here are applied post-hoc on these pre-trained feature representations (i.e., no additional training is performed for detection). Two broad families of post-hoc approaches are used in the experiments:

\begin{itemize}
    \item \textbf{Scoring-based methods:} Anomaly scores are computed from a single feature vector per sample using heuristics or calibrated transforms, for example Maximum Softmax Probability (MSP)~\cite{hendrycks2017baseline}, Energy-based detection~\cite{liu2020energy}, Virtual Logit Matching (ViM)~\cite{wang2022vim}, and ASH~\cite{djurisic2022ash}. These methods can be applied to spike-pattern features (penultimate and/or final-layer activations) or dense features (membrane potentials).
    \item \textbf{Distance-based methods:} Anomaly scores are derived from distances in feature space to in-distribution references (class prototypes or neighbors); examples include Nearest Class Mean (NCM)~\cite{huo2025ncm} and $k$-Nearest Neighbors (kNN)~\cite{sun2022knn}. For spiking networks, either binary spike patterns or membrane potentials can be used as feature vectors for computing these distances.
\end{itemize}

Fig.~\ref{fig:feature-extraction-pipeline-narrow} illustrates the experimental pipeline: feature extraction is separated from OoD scoring and all detectors are applied post-hoc on pre-trained feature vectors. Distance-based detectors (e.g., NCM, kNN) typically operate on penultimate-layer features to compute prototype or neighbor distances, while scoring-based detectors (e.g., MSP, Energy, ViM, ASH) use final- (and/or penultimate-) layer logits or transformed features to compute anomaly scores. When population coding is used, final-layer outputs are reshaped and aggregated per class before scoring.

\Figure[t!](topskip=0pt, botskip=0pt, midskip=0pt){figure1.eps}
{ \textbf{Input images are processed by the spiking backbone to produce features, optionally enhanced via population coding. The resulting features are then used for post-hoc or distance-based OoD detection. Dashed arrows indicate optional paths.}\label{fig:feature-extraction-pipeline-narrow}}

\subsection{Evaluation Metrics}
Performance of feature extraction and OoD detection is quantified using:
\begin{itemize}
    \item \textbf{AUROC:} Measures the probability that a randomly chosen ID sample receives a higher confidence score than a randomly chosen OoD sample. Higher AUROC indicates better discriminative capability.
    \item \textbf{FPR95:} Represents the fraction of OoD samples incorrectly classified as ID when the true positive rate for ID samples is 95\%. Lower FPR95 indicates more reliable OoD detection.
\end{itemize}


\section{Experimental results}
\label{sec:results}

\subsection{Inference-Time Computational Analysis}
As discussed in Sec.~\ref{sec:related}, CNN complexity is commonly reported using floating-point operation (FLOP) counts, typically expressed as the number of multiply--accumulate (MAC) operations that dominate convolutional inference. Architecture-level tools such as \texttt{ptflops}~\cite{ptflops} provide hardware-agnostic estimates of these costs. By contrast, spiking networks perform event-driven computation: synaptic operations occur only when spikes are emitted, so cost is more naturally expressed in synaptic operations (SynOps) or spike-triggered accumulations. In the spike-like models studied here, activations are binary and temporal integration is suppressed, which reduces arithmetic complexity.

To enable a practical comparison we compute inference-time operation counts for each spike-like network and its CNN counterpart while keeping the backbone architecture identical. In a CNN, each convolutional output element requires a MAC (multiplication plus accumulation). For spike-like networks processing static RGB inputs, MACs are required at the input/first layer (to process real-valued pixel intensities); in deeper layers binary activations allow multiplications to be replaced by accumulate-only (AC) operations. Normalization and some auxiliary layers still involve floating-point ops, but the dominant inference work in spike-like models shifts toward integer additions. Although the energy per operation depends on CMOS technology and implementation details~\cite{lemaire2023estimation}, aggregating MAC and AC counts provides a hardware-agnostic proxy for relative compute cost.

Table~\ref{tab:computation-comparison} summarizes these estimates for RGB inputs of size $32\times32$ (e.g., CIFAR). Values account for convolutional, normalization, and activation operations following the conventions in~\cite{ptflops, lemaire2023estimation}. Evaluated at a single time step, spike-like networks show large reductions in floating-point MAC counts, with most remaining compute expressed as ACs. As mentioned, 'spike-like' denotes single-step spiking models implemented on CNN backbones (CNN + LIF) and all reported spike-like operation counts are computed for single-step inference. The CNN column lists MAC counts for full-precision inference while the spike-like column separates ACs (deeper layers) and remaining MACs (e.g., first layer and auxiliary ops).

\begin{table}[tbp]
\caption{Inference-time computational complexity comparison between CNN and spike-like backbones (single-step inference, $T = 1$).}
\begin{center}
\small
\begin{tabular}{l p{2.75cm} p{2.75cm}}
\toprule
\textbf{Backbone} & \textbf{CNN (MACs)} & \textbf{Spike-like networks (ACs / MACs)} \\
\midrule
ConvNet  & 9.7~M    & 1.18~M   / 3.93~M \\
ResNet10 & 254.5~M  & 83.09~M  / 2.88~M \\
ResNet18 & 557.2~M  & 178.69~M / 3.72~M \\
\bottomrule
\end{tabular}
\label{tab:computation-comparison}
\end{center}
\end{table}

\subsection{Experimental setup}

Experiments follow OpenOOD protocols~\cite{yang2022openood}. Table~\ref{tab:ood-datasets} lists the near- and far-OoD datasets used for each in-distribution (ID) benchmark. All reported results show AUROC averaged over five independent training runs.

\begin{table}[tbp]
\caption{Near- and far-OoD datasets used for each in-distribution (ID) benchmark.}
\begin{center}
\small
\begin{tabular}{l p{2.75cm} p{2.75cm}}
\toprule
\textbf{ID dataset} & \textbf{Near-OoD} & \textbf{Far-OoD} \\
\midrule
CIFAR-10 & CIFAR-100, tImage200 & MNIST, SVHN, Textures, Places365 \\
SVHN & MNIST, EMNIST (digits) & CIFAR-10, Textures, Places365 \\
CIFAR-100 & CIFAR-10, tImage200 & MNIST, SVHN, Textures, Places365 \\
\bottomrule
\end{tabular}
\label{tab:ood-datasets}
\end{center}
\end{table}

To test the hypothesis that larger ID class counts benefit from larger population expansions, we evaluate a spectrum of expansion values tailored to each dataset: CIFAR-10: [1, 5, 10, 25, 50]; SVHN: [1, 5, 10, 25, 50, 100]; CIFAR-100: [25, 50, 100, 250, 500, 1000]. 

\subsection{Experiments with CIFAR10}

\subsubsection*{Penultimate features with nearest neighbors}

%\Figure[t!](topskip=0pt, botskip=0pt, midskip=0pt){figure2.eps}
%{ \textbf{AUROC for OoD detection on CIFAR-10 (in-distribution) as a function of population-coding expansion. Results are shown for spike-like models trained with MSE and CE objectives and a CNN baseline trained with CE, across three backbones and both near- and far-OoD datasets.}\label{fig:cifar10-experiment-1}}


\begin{table*}[tbp]
\caption{AUROC ($\uparrow$, upper sub-row) and FPR95 ($\downarrow$, lower sub-row) for OoD detection on CIFAR-10 (in-distribution) as a function of population-coding expansion. For AUROC rows, darker green indicates higher (better) values; for FPR95 rows, darker green indicates lower (better) values. Results are shown for spike-like models trained with MSE, CE and CL objectives and a CNN baseline trained with CE, across three backbones and both near- and far-OoD datasets.}
\begin{center}

\centering
\small
\renewcommand{\arraystretch}{1.18}
\setlength{\tabcolsep}{2pt}
\arrayrulecolor{gridgray}

\begin{tabular}{|L|C|C|C|C|C|G|C|C|C|C|C|}
\splitline
\rowcolor{headergray}
\multicolumn{1}{c}{} & \multicolumn{5}{c}{\textbf{Far-OoD}} & \multicolumn{1}{c}{\cellcolor{white}} & \multicolumn{5}{c}{\textbf{Near-OoD}} \\
\rowcolor{headergray}

\multicolumn{1}{c}{} & \multicolumn{1}{c}{\textbf{1}} & \multicolumn{1}{c}{\textbf{5}} & \multicolumn{1}{c}{\textbf{10}} & \multicolumn{1}{c}{\textbf{25}} & \multicolumn{1}{c}{\textbf{50}} & \multicolumn{1}{c}{\cellcolor{white}} & \multicolumn{1}{c}{\textbf{1}} & \multicolumn{1}{c}{\textbf{5}} & \multicolumn{1}{c}{\textbf{10}} & \multicolumn{1}{c}{\textbf{25}} & \multicolumn{1}{c}{\textbf{50}} \\


\splitline
\modelname{Conv-CE}      & \redcell{25.28} & \redcell{36.32} & \lowcell{50.24} & \lowcell{57.89} & \lowcell{61.78} & & \redcell{15.43} & \redcell{24.86} & \redcell{36.12} & \redcell{40.47} & \redcell{44.25} \\
\fprname & {---} & {---} & {---} & {---} & {---} & & {---} & {---} & {---} & {---} & {---} \\
\modelname{spike-Conv-MSE}     & \lowcell{68.34} & \lowcell{71.60} & \lowcell{69.53} & \lowcell{67.76} & \lowcell{68.39} & & \lowcell{54.18} & \lowcell{59.02} & \lowcell{59.55} & \lowcell{59.49} & \lowcell{53.79} \\
\fprname & {---} & {---} & {---} & {---} & {---} & & {---} & {---} & {---} & {---} & {---} \\
\modelname{spike-Conv-CE}      & \midcell{86.35} & \midcell{81.65} & \midcell{78.24} & \midcell{76.19} & \lowcell{73.71} & & \midcell{\textbf{84.13}} & \midcell{80.50} & \midcell{78.35} & \midcell{75.95} & \lowcell{73.28} \\
\fprname & {---} & {---} & {---} & {---} & {---} & & {---} & {---} & {---} & {---} & {---} \\
\modelname{spike-Conv-CL}      & \lowcell{62.52} & \lowcell{70.63} & \lowcell{70.52} & \lowcell{71.95} & \lowcell{72.70} & & \lowcell{54.98} & \lowcell{61.22} & \lowcell{60.56} & \lowcell{60.25} & \lowcell{60.09} \\
\fprname & {---} & {---} & {---} & {---} & {---} & & {---} & {---} & {---} & {---} & {---} \\

\splitline
\groupgap
\modelname{Resnet10-CE}  & \greenBcell{89.55} & \vhighcell{90.04} & \greenBcell{89.89} & \greenBcell{89.50} & \greenBcell{89.58} & & \greenAcell{88.37} & \greenAcell{88.60} & \greenAcell{88.52} & \greenAcell{88.05} & \greenAcell{88.16} \\
\fprname & {---} & {---} & {---} & {---} & {---} & & {---} & {---} & {---} & {---} & {---} \\
\modelname{spike-Resnet10-MSE} & \topcell{92.79} & \highcell{94.15} & \highcell{94.47} & \highcell{94.63} & \peakcell{\textbf{95.85}} & & \greenBcell{89.83} & \greenBcell{89.02} & \vhighcell{90.72} & \vhighcell{91.89} & \topcell{\textbf{92.83}} \\
\fprname & {---} & {---} & {---} & {---} & {---} & & {---} & {---} & {---} & {---} & {---} \\
\modelname{spike-Resnet10-CE}  & \midcell{86.12} & \lowcell{74.26} & \midcell{77.83} & \midcell{83.56} & \midcell{82.49} & & \midcell{81.98} & \lowcell{69.51} & \lowcell{72.88} & \midcell{79.60} & \midcell{78.18} \\
\fprname & {---} & {---} & {---} & {---} & {---} & & {---} & {---} & {---} & {---} & {---} \\
\modelname{spike-Resnet10-CL}  & \midcell{80.39} & \midcell{82.03} & \lowcell{65.89} & \lowcell{66.41} & \lowcell{66.46} & & \midcell{77.86} & \midcell{79.91} & \lowcell{62.89} & \lowcell{62.88} & \lowcell{63.38} \\
\fprname & {---} & {---} & {---} & {---} & {---} & & {---} & {---} & {---} & {---} & {---} \\

\splitline
\groupgap
\modelname{spike-Resnet18-CE}  & \topcell{93.85} & \highcell{95.46} & \highcell{95.48} & \highcell{95.34} & \peakcell{95.56} & & \vhighcell{92.16} & \topcell{93.36} & \topcell{93.46} & \topcell{93.35} & \topcell{93.56} \\
\fprname & {---} & {---} & {---} & {---} & {---} & & {---} & {---} & {---} & {---} & {---} \\
\modelname{spike-Resnet18-MSE} & \highcell{95.26} & \peakcell{\textbf{96.63}} & \peakcell{96.10} & \highcell{95.26} & \peakcell{95.91} & & \vhighcell{91.55} & \topcell{\textbf{93.91}} & \topcell{92.86} & \vhighcell{91.56} & \vhighcell{92.47} \\
\fprname & {---} & {---} & {---} & {---} & {---} & & {---} & {---} & {---} & {---} & {---} \\
\modelname{spike-Resnet18-CE}  & \greenAcell{88.33} & \greenAcell{88.57} & \vhighcell{90.41} & \vhighcell{90.83} & \vhighcell{91.82} & & \midcell{83.94} & \midcell{84.12} & \midcell{86.04} & \highmidcell{87.29} & \highmidcell{87.83} \\
\fprname & {---} & {---} & {---} & {---} & {---} & & {---} & {---} & {---} & {---} & {---} \\
\modelname{spike-Resnet18-CL}  & \midcell{76.09} & \greenAcell{87.25} & \lowcell{67.09} & \lowcell{64.94} & \lowcell{66.44} & & \lowcell{70.68} & \midcell{82.54} & \lowcell{62.20} & \lowcell{59.82} & \lowcell{61.05} \\
\fprname & {---} & {---} & {---} & {---} & {---} & & {---} & {---} & {---} & {---} & {---} \\
\splitline
\end{tabular}

\label{tab:cifar10-experiment-1}
\end{center}
\end{table*}

This experiment uses CIFAR-10 as the ID dataset and evaluates spike-like models trained with a single simulation time step. Feature vectors are extracted from the penultimate layer, and OoD detection is performed using a kNN distance-based detector. Each backbone model is evaluated across different population-coding expansions and the three training objectives described in Section~\ref{sec:methodology} (spike-count MSE, membrane-potential CE, and spike-count CL). For CNN baselines, dense penultimate-layer features are extracted from analogous backbones trained with CE to ensure a fair comparison.

Table~\ref{tab:cifar10-experiment-1} reports near-OoD and far-OoD AUROC for all three training objectives across three backbones and five population-coding expansions $(1, 5, 10, 25, 50)$. Cell background color encodes performance level: darker green indicates higher AUROC, yellow and orange denote intermediate values, and red marks the weakest scores. Bold entries indicate the best result per backbone within the Far-OoD and Near-OoD groups, respectively.

Empirically, MSE- and CL-trained models yield stronger kNN separation when membrane-potential vectors are used, whereas CE-trained models perform better with spike-pattern features; the results in Table~\ref{tab:cifar10-experiment-1} therefore reflect this pairing throughout. Nearest-neighbor distances are computed directly on these features without additional detector training. Although population coding is applied at the output layer, its effect on OoD detection is assessed through its influence on the learned penultimate representations.

Spike-like models trained with the MSE objective consistently yield penultimate-layer membrane-potential features that achieve strong kNN-based OoD separation, while spike-like models trained with the CE objective produce competitive penultimate-layer spike-pattern features. Models trained with the CL (count loss) objective produce substantially weaker OoD representations across all backbones and population-coding expansions.

When population coding is enabled, spike-like models trained with MSE and evaluated using penultimate-layer features achieve AUROC values comparable to, and in some cases exceeding, those of CNN baselines. These results support the hypothesis that single-step spiking models can produce discriminative OoD representations when combined with population coding. Moreover, population coding improves OoD performance across all backbones, with the larger gains observed for residual architectures. Although population coding is applied only at the final layer, its effect propagates to the penultimate layer through backpropagation: each penultimate neuron receives gradient contributions from $E$ output neurons per class rather than one, so the training loss enforces class separation through more gradient paths simultaneously, promoting more class-discriminative penultimate representations. As a result, features extracted before the population-coded layer exhibit improved intra-class compactness and inter-class separation, which directly benefits distance-based OoD detection.


\subsubsection*{Final-layer features with scoring-based post-hoc methods}

%\Figure[t!](topskip=0pt, botskip=0pt, midskip=0pt){figure3.eps}
%{\textbf{AUROC for scoring-based post-hoc OoD detection on CIFAR-10 (in-distribution) as a function of population-coding expansion. Results are shown for spike-like residual models trained with MSE, across near- and far-OoD datasets; solid and dashed curves correspond to membrane-potential and spike-pattern features, respectively. Horizontal reference lines indicate CNN baseline AUROC values for each method.}\label{fig:cifar10-experiment-2}}


\begin{table*}[tbp]
\caption{AUROC for scoring-based post-hoc OoD detection on CIFAR-10 (in-distribution) as a function of population-coding expansion. Results are shown for spike-like residual models trained with the MSE objective, across near- and far-OoD datasets. Only membrane-potential features are reported; spike-pattern features yielded substantially lower AUROC across all evaluated methods. The \textbf{CNN} column reports the single baseline value for the corresponding conventional CNN backbone (no population coding).}
\begin{center}

\centering
\small
\renewcommand{\arraystretch}{1.18}
\setlength{\tabcolsep}{2pt}
\arrayrulecolor{gridgray}

% Column layout: L | C×5 (spike E=1..50) | C (CNN ref) | G (gap) | C×5 (spike E=1..50) | C (CNN ref)
% CNN column carries the single baseline value; spike rows leave it as ---
\begin{tabular}{|L|C|C|C|C|C|C|G|C|C|C|C|C|C|}
\cline{1-7}\cline{9-14}
\rowcolor{headergray}
\multicolumn{1}{c}{} & \multicolumn{6}{c}{\textbf{Far-OoD}} & \multicolumn{1}{c}{\cellcolor{white}} & \multicolumn{6}{c}{\textbf{Near-OoD}} \\
\rowcolor{headergray}
\multicolumn{1}{c}{} & \multicolumn{1}{c}{\textbf{1}} & \multicolumn{1}{c}{\textbf{5}} & \multicolumn{1}{c}{\textbf{10}} & \multicolumn{1}{c}{\textbf{25}} & \multicolumn{1}{c}{\textbf{50}} & \multicolumn{1}{c}{\textbf{CNN}} & \multicolumn{1}{c}{\cellcolor{white}} & \multicolumn{1}{c}{\textbf{1}} & \multicolumn{1}{c}{\textbf{5}} & \multicolumn{1}{c}{\textbf{10}} & \multicolumn{1}{c}{\textbf{25}} & \multicolumn{1}{c}{\textbf{50}} & \multicolumn{1}{c}{\textbf{CNN}} \\

\cline{1-7}\cline{9-14}
\groupgap
\modelname{spike-Resnet10 ASH}    & \lowcell{69.94} & \midcell{86.97} & \greenAcell{88.12} & \vhighcell{90.37} & \greenBcell{89.22} & 89.01 & & \lowcell{56.16} & \midcell{80.85} & \midcell{82.99} & \midcell{85.39} & \midcell{81.47} & 83.72 \\
\modelname{spike-Resnet10 MSP}    & \lowcell{73.07} & \midcell{85.83} & \highmidcell{87.66} & \greenBcell{89.39} & \greenBcell{89.89} & 87.66 & & \lowcell{74.13} & \midcell{79.92} & \midcell{82.69} & \midcell{84.59} & \midcell{83.26} & 84.63 \\
\modelname{spike-Resnet10 ODIN}   & \lowcell{73.01} & \midcell{85.70} & \greenAcell{88.24} & \vhighcell{90.31} & \vhighcell{90.18} & 87.73 & & \lowcell{74.08} & \midcell{79.79} & \midcell{83.16} & \midcell{85.60} & \midcell{83.39} & 82.36 \\
\modelname{spike-Resnet10 Energy} & \redcell{44.45} & \highmidcell{87.52} & \midcell{84.10} & \midcell{86.22} & \midcell{78.60} & 89.01 & & \redcell{31.14} & \midcell{81.63} & \midcell{80.46} & \midcell{82.66} & \lowcell{71.05} & 83.70 \\
\modelname{spike-Resnet10 MLS}    & \midcell{82.89} & \midcell{86.60} & \highmidcell{87.52} & \greenBcell{89.56} & \greenBcell{89.38} & 89.92 & & \lowcell{75.09} & \midcell{80.53} & \midcell{82.52} & \midcell{84.61} & \midcell{81.93} & 83.81 \\
\modelname{spike-Resnet10 VIM}    & \midcell{85.34} & \peakcell{\textbf{96.77}} & \peakcell{95.55} & \peakcell{95.55} & \topcell{93.31} & 94.44 & & \midcell{81.05} & \topcell{\textbf{92.75}} & \vhighcell{91.13} & \vhighcell{91.97} & \midcell{86.68} & 91.95 \\

\cline{1-7}\cline{9-14}
\groupgap
% TODO: fill CNN column — replace --- with \<color>cell{value} once available
\modelname{spike-Resnet18 ASH}    & \lowcell{73.21} & \midcell{85.70} & \topcell{92.10} & \topcell{92.92} & \topcell{93.90} & 92.87 & & \lowcell{61.05} & \midcell{86.91} & \greenAcell{88.45} & \greenBcell{89.60} & \greenBcell{89.56} & 89.40 \\
\modelname{spike-Resnet18 MSP}    & \lowcell{75.92} & \vhighcell{91.75} & \topcell{92.63} & \topcell{92.67} & \topcell{92.38} & 91.74 & & \lowcell{70.55} & \midcell{86.48} & \greenAcell{88.14} & \greenAcell{88.00} & \highmidcell{87.15} & 89.43 \\
\modelname{spike-Resnet18 ODIN}   & \lowcell{75.89} & \vhighcell{91.75} & \topcell{92.83} & \topcell{93.65} & \topcell{93.42} & 91.93 & & \lowcell{70.51} & \midcell{86.44} & \greenAcell{88.38} & \greenBcell{89.28} & \greenAcell{88.47} & 88.41 \\
\modelname{spike-Resnet18 Energy} & \redcell{47.93} & \vhighcell{91.18} & \vhighcell{91.84} & \topcell{93.88} & \highcell{94.07} & 92.86 & & \redcell{40.29} & \midcell{86.81} & \greenAcell{88.10} & \greenBcell{89.62} & \greenBcell{89.61} & 89.39 \\
\modelname{spike-Resnet18 MLS}    & \midcell{83.63} & \vhighcell{91.84} & \topcell{92.60} & \topcell{93.48} & \topcell{93.43} & 92.83 & & \lowcell{75.07} & \midcell{86.64} & \greenAcell{88.09} & \greenBcell{89.09} & \greenAcell{88.54} & 89.41 \\
\modelname{spike-Resnet18 VIM}    & \greenAcell{88.39} & \topcell{92.00} & \topcell{93.01} & \topcell{93.87} & \peakcell{\textbf{95.57}} & 93.36 & & \midcell{82.20} & \highmidcell{87.99} & \greenBcell{89.39} & \vhighcell{90.07} & \vhighcell{\textbf{91.69}} & 91.90 \\
\cline{1-7}\cline{9-14}
\end{tabular}

\label{tab:cifar10-experiment-2}
\end{center}
\end{table*}


This experiment evaluates scoring-based OoD detection methods applied to membrane-potential features extracted from spike-like residual networks trained with the MSE objective and operating at a single simulation time step. The evaluated methods include Maximum Softmax Probability (MSP), Energy, Maximum Logit Score (MLS), ASH, and Virtual Logit Matching (ViM). All methods are applied without additional detector training. In contrast to the previous experiment, where distance-based kNN detection was performed on penultimate-layer features, here post-hoc anomaly scores are computed from final-layer representations (with exception of ViM which uses both penultimate and final layer features). Population-coding expansions are varied to assess whether distributed output representations influence the separability of final-layer features and the effectiveness of scoring-based OoD detection.

For the spike-ConvNet backbone, the resulting anomaly scores are uninformative regardless of the training objective, and its results are therefore omitted. Furthermore, spike-pattern features yield substantially lower AUROC than membrane-potential features across all evaluated methods, so only membrane-potential feature results are reported.

Table~\ref{tab:cifar10-experiment-2} summarizes results for the spike-ResNet10 and spike-ResNet18 backbones. Several consistent observations emerge. First, population coding consistently improves OoD detection performance across multiple post-hoc methods: demonstrating that distributed output representations benefit final-layer scoring-based detection as well as penultimate-layer distance-based detection. Second, ViM, which leverages both penultimate- and final-layer information, generally achieves the strongest performance, reaching peak AUROC values that substantially exceed those of the other methods. In some cases, the anomaly score  is especially sensitive to the absence of population coding, exhibiting near-random performance at $E=1$ for both backbones while recovering strongly at larger expansion values.

These results suggest that single-step spike-like models trained with the MSE objective can enable effective application of standard post-hoc OoD scoring methods, particularly when dense membrane-potential features are combined with population coding.



\subsubsection*{Effect of multi-step training on single-step OoD Detection}

%\begin{table}[tbp]
%\caption{Effect of multi-step training on kNN and penultimate features. The given numbers is auroc for Near- and far-ood datasets.}
%\begin{center}
%\begin{tabular}{l| p{1cm} p{1cm} p{1cm} p{1cm} p{1cm}}
%    & \multicolumn{5}{l}{spike-ResNet18} \\
%    & 1             & 5             & 10            & 25            & 50            \\ \hline
%T=1 & 91.55 / 95.26 & 93.91 / 96.63 & 92.86 / 96.10 & 91.56 / 95.26 & 92.47 / 95.91 \\
%T=2 & 53.86 / 55.28 & 76.90 / 81.24 & 84.12 / 90.31 & 96.70 / 98.63 & 97.25 / 98.52 \\
%T=4 & 51.60 / 53.64 & 79.61 / 86.76 & 88.02 / 93.70 & 95.84 / 98.01 & 96.33 / 98.15 \\
%T=8 & 84.13 / 99.62 & 94.38 / 98.21 & 96.19 / 97.73 & 96.80 / 98.32 & 95.23 / 97.89    
%\end{tabular}
%\label{tab:cifar10-experiment-3}
%\end{center}
%\end{table}
%
%
%This experiment introduces ablation setup based on the previous setup and employing multiple simulation time steps during training, while keeping inference restricted to a single time step. This way, training is leaned towards the SNN approach. It is investigated whether OoD detection performance can benefit from multi-step training while preserving the energy efficiency of single-step inference.
%
%Specifically, the spike-ResNet18 backbone is tested to examine the effect of multi-step training on kNN based OoD (like in the first experiment) and the best scoring method from the second experiment. The underlying hypothesis is that temporal integration during training may refine synaptic connections and internal representations, leading to more discriminative feature vectors. If multi-step training enables the network to better capture the structure of Id data, it is expected that Id and OoD inputs will produce more separable spike patterns or membrane-potential features. This experiment also includes population coding, allowing us to assess whether the combination of temporal dynamics and population-based output representations provides complementary benefits for OoD detection.
%
%The Table~\ref{tab:cifar10-experiment-3} gives comparison of AUROC for spike-ResNet18 for different number of training time steps. It can be noted that increasing the training time steps increases the AUROC value for larger expansion values of population coding. 


\begin{table*}[tbp]
\caption{AUROC for kNN-based OoD detection on penultimate-layer membrane-potential features of spike-ResNet18 trained with the MSE objective, as a function of training time steps $T$ and population-coding expansion $E$. Inference is restricted to a single time step in all cases. Results are averaged over five independent runs.}
\begin{center}
\centering
\small
\renewcommand{\arraystretch}{1.18}
\setlength{\tabcolsep}{2pt}
\arrayrulecolor{gridgray}

\begin{tabular}{|L|C|C|C|C|C|G|C|C|C|C|C|}
\splitline
\rowcolor{headergray}
\multicolumn{1}{c}{} & \multicolumn{5}{c}{\textbf{Far-OoD}} & \multicolumn{1}{c}{\cellcolor{white}} & \multicolumn{5}{c}{\textbf{Near-OoD}} \\
\rowcolor{headergray}
\multicolumn{1}{c}{} & \multicolumn{1}{c}{\textbf{1}} & \multicolumn{1}{c}{\textbf{5}} & \multicolumn{1}{c}{\textbf{10}} & \multicolumn{1}{c}{\textbf{25}} & \multicolumn{1}{c}{\textbf{50}} & \multicolumn{1}{c}{\cellcolor{white}} & \multicolumn{1}{c}{\textbf{1}} & \multicolumn{1}{c}{\textbf{5}} & \multicolumn{1}{c}{\textbf{10}} & \multicolumn{1}{c}{\textbf{25}} & \multicolumn{1}{c}{\textbf{50}} \\

\splitline
kNN $T=1$ & \highcell{95.26} & \peakcell{96.63} & \peakcell{96.10} & \highcell{95.26} & \peakcell{95.91} & & \vhighcell{91.55} & \topcell{93.91} & \topcell{92.86} & \vhighcell{91.56} & \topcell{92.47} \\
kNN $T=2$ & \lowcell{55.28} & \midcell{81.24} & \vhighcell{90.31} & \peakcell{98.63} & \peakcell{98.52} & & \lowcell{53.86} & \midcell{76.90} & \midcell{84.12} & \peakcell{96.70} & \peakcell{97.25} \\
kNN $T=4$ & \lowcell{53.64} & \midcell{86.76} & \topcell{93.70} & \peakcell{98.01} & \peakcell{98.15} & & \lowcell{51.60} & \midcell{79.61} & \greenAcell{88.02} & \peakcell{95.84} & \peakcell{96.33} \\
kNN $T=8$ & \peakcell{99.62} & \peakcell{98.21} & \peakcell{97.73} & \peakcell{98.32} & \peakcell{97.89} & & \midcell{84.13} & \highcell{94.38} & \peakcell{96.19} & \peakcell{96.80} & \highcell{95.23} \\
\splitline
\end{tabular}

\label{tab:cifar10-experiment-3}
\end{center}
\end{table*}

\begin{table*}[tbp]
\caption{AUROC for ViM-based OoD detection on spike-ResNet18 trained with the MSE objective, as a function of training time steps $T$ and population-coding expansion $E$. Inference is restricted to a single time step in all cases. Results are averaged over five independent runs. }
\begin{center}
\centering
\small
\renewcommand{\arraystretch}{1.18}
\setlength{\tabcolsep}{2pt}
\arrayrulecolor{gridgray}

\begin{tabular}{|L|C|C|C|C|C|G|C|C|C|C|C|}
\splitline
\rowcolor{headergray}
\multicolumn{1}{c}{} & \multicolumn{5}{c}{\textbf{Far-OoD}} & \multicolumn{1}{c}{\cellcolor{white}} & \multicolumn{5}{c}{\textbf{Near-OoD}} \\
\rowcolor{headergray}
\multicolumn{1}{c}{} & \multicolumn{1}{c}{\textbf{1}} & \multicolumn{1}{c}{\textbf{5}} & \multicolumn{1}{c}{\textbf{10}} & \multicolumn{1}{c}{\textbf{25}} & \multicolumn{1}{c}{\textbf{50}} & \multicolumn{1}{c}{\cellcolor{white}} & \multicolumn{1}{c}{\textbf{1}} & \multicolumn{1}{c}{\textbf{5}} & \multicolumn{1}{c}{\textbf{10}} & \multicolumn{1}{c}{\textbf{25}} & \multicolumn{1}{c}{\textbf{50}} \\

\splitline
ViM $T=1$ & \greenAcell{88.39} & \topcell{92.00} & \peakcell{96.10} & \highcell{95.26} & \peakcell{95.57} & & \midcell{82.20} & \highmidcell{87.99} & \topcell{93.01} & \topcell{93.87} & \vhighcell{91.69} \\
ViM $T=2$ & \lowcell{54.08} & \midcell{82.03} & \greenBcell{89.91} & \highcell{95.46} & \highcell{95.46} & & \lowcell{53.49} & \lowcell{75.93} & \midcell{84.15} & \vhighcell{91.43} & \topcell{92.53} \\
ViM $T=4$ & \lowcell{53.11} & \midcell{85.14} & \vhighcell{90.17} & \topcell{92.45} & \topcell{93.36} & & \lowcell{51.24} & \midcell{77.46} & \midcell{83.67} & \midcell{86.23} & \greenAcell{88.79} \\
ViM $T=8$ & \midcell{81.78} & {---} & \vhighcell{91.10} & \highcell{94.85} & \peakcell{98.65} & & \lowcell{67.11} & {---} & \midcell{85.77} & \topcell{92.80} & \highcell{94.59} \\
\splitline
\end{tabular}

\label{tab:cifar10-experiment-3-vim}
\end{center}
\end{table*}

This experiment extends the first CIFAR-10 experiment into an ablation study: the spike-ResNet18 backbone is trained with the MSE objective across four training-step budgets $T \in \{1, 2, 4, 8\}$, while inference is kept at a single time step in all cases. Restricting inference to $T=1$ preserves the low-compute, energy-efficient profile of spike-like models regardless of how many steps are used during training. The kNN distance-based detector applied to penultimate-layer membrane-potential features is retained throughout, isolating the effect of training depth from other variables. Table~\ref{tab:cifar10-experiment-3} and Table~\ref{tab:cifar10-experiment-3-vim} report results for kNN and ViM detection, respectively.

Two distinct regimes emerge from the color patterns in the tables. Without population coding ($E=1$), increasing the number of training steps does not consistently improve OoD performance and in several cases substantially degrades it. This indicates that temporal integration during training restructures the internal representations in a way that is poorly suited to single-step feature extraction when the output layer is unexpanded. With sufficient population coding ($E \geq 25$), the picture changes markedly: multi-step training then yields clearly stronger OoD separation than single-step training, with the improvement visible as a shift from lighter to darker greens across the corresponding columns. The benefit of multi-step training therefore depends on the presence of an adequately expanded output layer to relay the richer gradient signal back to the penultimate layer.

Population coding plays a dual role here. In addition to amplifying the effect of multi-step training, it independently improves OoD performance even at $T=1$, as already observed in the preceding experiments. The two mechanisms are therefore complementary: population coding stabilizes and enhances the penultimate-layer feature structure, while multi-step training further refines it when population coding provides sufficient training signal.

A further observation concerns the relative merit of detection strategies. The kNN detector applied to penultimate-layer membrane-potential features achieves performance comparable to, and in several configurations exceeding, the best scoring-based method (ViM) evaluated on final-layer features. This is notable because kNN operates directly on raw penultimate features without any calibration or transformation, suggesting that well-structured penultimate representations alone can be sufficient for competitive OoD detection.
    
 
\begin{table}[tbp]
\caption{Top-5 entries on the OpenOOD~v1.5 CIFAR-10 leaderboard (ResNet-18,
$\approx$95\% ID accuracy), ranked by Near-OoD AUROC. Near-OoD: CIFAR-100,
TinyImageNet. Far-OoD: MNIST, SVHN, Textures, Places365. $^\dagger$Uses
outlier training data (TIN-597).}
\begin{center}
\footnotesize
\begin{tabular}{|c|p{2.7cm}|p{2.4cm}|c||p{2.7cm}|p{2.4cm}|c|}
\hline
\multirow{2}{*}{\textbf{Rank}}
  & \multicolumn{3}{c||}{\textbf{Near-OoD AUROC}}
  & \multicolumn{3}{c|}{\textbf{Far-OoD AUROC}} \\
\cline{2-7}
  & \textbf{Training} & \textbf{Postprocessor} & \textbf{AUROC}
  & \textbf{Training} & \textbf{Postprocessor} & \textbf{AUROC} \\
\hline
1 & RotPred + PixMix & RotPred         & 94.86
  & RotPred + PixMix & RotPred         & 98.18 \\
2 & OE$^\dagger$     & MSP             & 94.82
  & LogitNorm        & MSP             & 96.74 \\
3 & RotPred          & RotPred         & 92.68
  & RotPred          & RotPred         & 96.62 \\
4 & LogitNorm        & MSP             & 92.33
  & OE$^\dagger$     & MSP             & 96.00 \\
5 & MCD$^\dagger$    & MCD             & 91.03
  & G-ODIN           & G-ODIN          & 95.51 \\
\hline
\end{tabular}
\label{tab:openood-cifar10}
\end{center}
\end{table}
    
\begin{table}[tbp]
\caption{Top-5 entries on the OpenOOD~v1.5 CIFAR-100 leaderboard (ResNet-18,
$\approx$77\% ID accuracy), ranked by Near-OoD AUROC. Near-OoD: CIFAR-10,
TinyImageNet. Far-OoD: MNIST, SVHN, Textures, Places365. $^\dagger$Uses
outlier training data (TIN-597). $^\ddagger$AUROC unavailable due to
leaderboard pagination.}
\begin{center}
\footnotesize
\begin{tabular}{|c|p{2.7cm}|p{2.4cm}|c||p{2.7cm}|p{2.4cm}|c|}
\hline
\multirow{2}{*}{\textbf{Rank}}
  & \multicolumn{3}{c||}{\textbf{Near-OoD AUROC}}
  & \multicolumn{3}{c|}{\textbf{Far-OoD AUROC}} \\
\cline{2-7}
  & \textbf{Training} & \textbf{Postprocessor} & \textbf{AUROC}
  & \textbf{Training} & \textbf{Postprocessor} & \textbf{AUROC} \\
\hline
1 & OE$^\dagger$   & MSP         & 88.30
  & RotPred        & RotPred     & 88.40 \\
2 & OAL            & EBO         & 82.03
  & G-ODIN         & G-ODIN      & 85.67 \\
3 & CrossEntropy   & WeiPer+KLD  & 81.37
  & OAL            & EBO         & 84.03 \\
4 & CrossEntropy   & GEN         & $-^\ddagger$
  & CrossEntropy   & RMDS        & 82.92 \\
5 & CrossEntropy   & MLS         & 81.05
  & CrossEntropy   & KNN         & 82.40 \\
\hline
\end{tabular}
\label{tab:openood-cifar100}
\end{center}
\end{table}

\subsection{Experiment 4: SVHN and CIFAR100 with single-step model}
Just keep it as a proof of a concept, avoiding detailed graphic representations.


TODO: Consider where to input the FPR95 in the tables 

\section{Conclusions}
\label{sec:conclusions}

This paper proposed spike-like models for energy-efficient out-of-distribution detection. The proposed solution is based on standard convolutional backbones with leaky integrate-and-fire activations constrained to single-step inference. By replacing dense floating-point activations with binary spiking activations, the approach substantially reduces arithmetic complexity relative to conventional CNN-based OoD methods, while retaining compatibility with standard post-hoc detection pipelines. The primary motivation is to achieve OoD performance comparable to or better than CNN baselines reported in OpenOOD, at reduced computational cost.

Experiments on CIFAR-10 systematically examined three training objectives (MSE, CE, and CL), three backbone architectures (spike-ConvNet, spike-ResNet10, spike-ResNet18), and a range of population-coding expansion values. The first experimental stage evaluated distance-based detection: kNN applied to penultimate-layer membrane-potential features. Spike-like models trained with the MSE objective consistently produce well-structured membrane-potential representations, and with sufficient population coding their kNN AUROC matches or exceeds that of conventional CNN baselines. The CL objective, by contrast, yields substantially weaker feature representations across all settings.

The second stage evaluated scoring-based post-hoc methods on final-layer membrane-potential features, also using single-step training. These methods — including MSP, Energy, ODIN, MLS, ASH, and ViM — likewise benefit from population coding, with AUROC values improving progressively as the expansion increases. ViM, which incorporates both penultimate- and final-layer information, achieves the highest overall performance. Notably, the simpler kNN detector on raw penultimate features remains competitive with the best scoring-based results, without requiring any post-hoc calibration or transformation.

The third stage extended the ablation to multi-step training while retaining single-step inference, preserving the low-power deployment profile. This analysis demonstrated that additional training time steps can further improve OoD separability, but only when population coding provides a sufficiently expanded output layer. Without adequate expansion, multi-step training disrupts rather than improves the penultimate-layer feature structure. With it, training depth and population coding interact synergistically, jointly shaping more discriminative representations.

Throughout all experiments, population coding emerges as the single most consistent performance enabler: it improves distance-based and scoring-based detection independently, and it is a prerequisite for multi-step training to be beneficial. These findings support the broader hypothesis that population coding can partially substitute for temporal training depth, making it a practical tool for obtaining competitive, energy-efficient OoD detection from spike-like models operating under single-step inference constraints.



\section*{Acknowledgment}
The preferred spelling of the word ``acknowledgment'' in American English is
without an ``e'' after the ``g.'' Use the singular heading even if you have
many acknowledgments. Avoid expressions such as ``One of us (S.B.A.) would
like to thank $\ldots$ .'' Instead, write ``F. A. Author thanks $\ldots$ .'' In most
cases, sponsor and financial support acknowledgments are placed in the
unnumbered footnote on the first page, not here.

\bibliographystyle{IEEEtran}
\bibliography{references} 


\begin{IEEEbiography}[{\includegraphics[width=1in,height=1.25in,clip,keepaspectratio]{author1.png}}]{First A. Author} received the B.S. and M.S. degrees in aerospace engineering from
the University of Virginia, Charlottesville, in 2001 and the Ph.D. degree in
mechanical engineering from Drexel University, Philadelphia, PA, in 2008.

From 2001 to 2004, he was a Research Assistant with the Princeton Plasma
Physics Laboratory. Since 2009, he has been an Assistant Professor with the
Mechanical Engineering Department, Texas A{\&}M University, College Station.
He is the author of three books, more than 150 articles, and more than 70
inventions. His research interests include high-pressure and high-density
nonthermal plasma discharge processes and applications, microscale plasma
discharges, discharges in liquids, spectroscopic diagnostics, plasma
propulsion, and innovation plasma applications. He is an Associate Editor of
the journal \emph{Earth, Moon, Planets}, and holds two patents.

Dr. Author was a recipient of the International Association of Geomagnetism
and Aeronomy Young Scientist Award for Excellence in 2008, and the IEEE
Electromagnetic Compatibility Society Best Symposium Paper Award in 2011.
\end{IEEEbiography}


\begin{IEEEbiography}[{\includegraphics[width=1in,height=1.25in,clip,keepaspectratio]{author2.png}}]{Second B. Author} (M'76--SM'81--F'87) and all authors may include
biographies. Biographies are often not included in conference-related
papers. This author became a Member (M) of IEEE in 1976, a Senior
Member (SM) in 1981, and a Fellow (F) in 1987. The first paragraph may
contain a place and/or date of birth (list place, then date). Next,
the author's educational background is listed. The degrees should be
listed with type of degree in what field, which institution, city,
state, and country, and year the degree was earned. The author's major
field of study should be lower-cased.

The second paragraph uses the pronoun of the person (he or she) and not the
author's last name. It lists military and work experience, including summer
and fellowship jobs. Job titles are capitalized. The current job must have a
location; previous positions may be listed
without one. Information concerning previous publications may be included.
Try not to list more than three books or published articles. The format for
listing publishers of a book within the biography is: title of book
(publisher name, year) similar to a reference. Current and previous research
interests end the paragraph.

The third paragraph begins with the author's
title and last name (e.g., Dr.\ Smith, Prof.\ Jones, Mr.\ Kajor, Ms.\ Hunter).
List any memberships in professional societies other than the IEEE. Finally,
list any awards and work for IEEE committees and publications. If a
photograph is provided, it should be of good quality, and
professional-looking. Following are two examples of an author's biography.
\end{IEEEbiography}

\newpage

%If you do not have or do not want to include a photo, you can use IEEEbiographynophoto as shown below:

\begin{IEEEbiographynophoto}{Third C. Author, Jr.} (M'87) received the B.S. degree in mechanical
engineering from National Chung Cheng University, Chiayi, Taiwan, in 2004
and the M.S. degree in mechanical engineering from National Tsing Hua
University, Hsinchu, Taiwan, in 2006. He is currently pursuing the Ph.D.
degree in mechanical engineering at Texas A{\&}M University, College
Station, TX, USA.

From 2008 to 2009, he was a Research Assistant with the Institute of
Physics, Academia Sinica, Tapei, Taiwan. His research interest includes the
development of surface processing and biological/medical treatment
techniques using nonthermal atmospheric pressure plasmas, fundamental study
of plasma sources, and fabrication of micro- or nanostructured surfaces.

Mr. Author's awards and honors include the Frew Fellowship (Australian
Academy of Science), the I. I. Rabi Prize (APS), the European Frequency and
Time Forum Award, the Carl Zeiss Research Award, the William F. Meggers
Award and the Adolph Lomb Medal (OSA).
\end{IEEEbiographynophoto}

\EOD

\end{document}
